% متن ارائه‌ی شفاهی آزمایش ۴ — نسخه‌ی موقّت، فقط برای خواندن از روی آن
\documentclass[12pt,a4paper]{article}
\usepackage{../latex/dsdlab-report}

\linespread{1.6}
\setlength{\parskip}{0.5em}
\setlength{\parindent}{0pt}

\definecolor{cuecolor}{HTML}{9A3412}
\newcommand{\cue}[2]{%
  \par\addvspace{1.1em}%
  \noindent{\small\bfseries\color{cuecolor}\lr{#1}\quad\textbullet\quad #2}%
  \par\nobreak\addvspace{0.3em}%
}

\begin{document}
\thispagestyle{empty}

\begin{center}
  {\Large\bfseries متن ارائه‌ی آزمایش ۴ --- توصیف رفتاری}\\[0.3em]
  {\small حدود ۵ دقیقه \quad|\quad پوریا عرفانزاده، فاطیما تیمارچی، آیه صابری}
\end{center}
\vspace{0.6em}
\hrule
\vspace{0.4em}

\large

\cue{0:00}{شروع}
سلام. من پوریا عرفانزاده هستم، با فاطیما تیمارچی و آیه صابری. آزمایش چهارم ما توصیف رفتاری بود و چیزی که ساختیم یه استک سخت‌افزاری با عمق ۸ و پهنای ۴ بیته.

\cue{0:20}{صورت مسئله}
استک همون ساختار \lr{LIFO} آشناست: هرچی آخر گذاشتی، اوّل برمی‌داری. ورودی‌هاش کلاک، ریست، داده، و دو تا فرمان \lr{Push} و \lr{Pop}ه؛ خروجی‌هاش داده‌ی خروجی و دو تا فلگ \lr{Full} و \lr{Empty}.

نکته‌ی این آزمایش این بود که برعکس آزمایش سوم که فقط \lr{dataflow} مجاز بود، اینجا فقط توصیف رفتاری مجاز بود --- یعنی \lr{always} و \lr{case} و انتساب \lr{non-blocking}. ما رفتار مطلوب رو می‌نویسیم و سنتزکننده خودش تصمیم می‌گیره چه سخت‌افزاری بسازه.

\cue{0:55}{دو تا نکته‌ی ریز ولی مهم}
اوّل اینکه فلگ \lr{Empty} تو دستور کار فعّال‌صفره. صریح نوشته \lr{Empty=0} یعنی استک خالیه. یعنی برخلاف انتظار، وقتی \lr{Empty} بالاست یعنی استک خالی \emph{نیست}. این رو تو کل کد و تست‌بنچ رعایت کردیم و شکل موج‌ها رو هم باید با همین قرارداد خوند.

دوم اینکه دستور کار چند تا حالت مرزی رو اصلاً تعریف نکرده بود و خودمون باید تصمیم می‌گرفتیم. \lr{Push} روی استک پر رو نادیده می‌گیریم که داده‌ی معتبر بازنویسی نشه. \lr{Pop} از استک خالی رو هم نادیده می‌گیریم و \lr{Data\_Out} مقدار قبلیش رو نگه می‌داره --- صفر کردنش کار درستی نبود، چون با یه داده‌ی واقعیِ صفر قاطی می‌شد. و \lr{Push} و \lr{Pop} هم‌زمان: هیچ کاری نمی‌کنیم. این ساده‌ترین و قابل‌تست‌ترین انتخاب بود و ما رو از تعریف یه اولویّت دلبخواهی نجات داد.

\cue{1:45}{اشاره‌گر و پهنای آن}
هسته‌ی طرح یه رجیستره به اسم \lr{Stack\_Pointer} که تعداد عناصر معتبر رو نگه می‌داره. چون از صفر می‌شماره، همین عدد هم‌زمان آدرس اوّلین خانه‌ی خالی هم هست؛ پس \lr{Push} می‌نویسه روی خودِ اشاره‌گر و بعد یکی زیاد می‌کنه، و \lr{Pop} از یکی کمتر می‌خونه و بعد یکی کم می‌کنه. متقارن و ساده.

یه نکته‌ی ظریف تو پهنای اشاره‌گره. ممکنه فکر کنید برای عمق ۸، سه بیت کافیه؛ ولی کافی نیست. اشاره‌گر باید ۹ تا مقدار مختلف رو نشون بده، از صفر تا ۸. اگه سه بیت می‌گرفتیم، مقدار ۸ که دقیقاً همون شرط پر بودنه اصلاً نمایش‌پذیر نبود و شمارنده به‌جاش برمی‌گشت رو صفر. پس فرمولش لگاریتم دوی \emph{عمق به‌علاوه‌ی یک} می‌شه، که برای ما چهار بیته.

هر دو فلگ هم فقط از روی همین اشاره‌گر و به‌صورت ترکیبی ساخته می‌شن. اینکه هر دو از یه منبع واحد بیان، امکان ناسازگاری بین وضعیّت واقعی استک و چیزی که گزارش می‌شه رو از بین می‌بره.

\cue{2:35}{تست‌بنچ}
تست‌بنچ خودآزماست و هفت مرحله داره: ریست، \lr{Pop} از خالی، پر کردن کامل، \lr{Push} روی پر، برداشت کامل با بررسی ترتیب \lr{LIFO}، دوباره \lr{Pop} از خالی، و آخرش فرمان هم‌زمان.

یه نکته‌ی مهم تو مرحله‌ی آخر هست: اونجا علاوه بر خروجی‌ها، خودِ اشاره‌گر رو هم مستقیم چک می‌کنیم. دلیلش اینه که اگه فقط به خروجی‌ها نگاه می‌کردیم، یه مدار غلط که هم \lr{Push} و هم \lr{Pop} رو انجام می‌ده هم از تست سربلند بیرون میومد --- چون اثر دو تا عمل روی اشاره‌گر همدیگه رو خنثی می‌کنه. سنجش مستقیم اشاره‌گر این راه فرار رو می‌بنده.

\cue{3:25}{پارامتری بودن}
طرح با دو تا پارامتر \lr{DEPTH} و \lr{WIDTH} نوشته شده و هیچ‌جای کد به عددهای ۸ و ۴ وابستگی سخت نداره. برای اینکه این فقط ادّعا نمونه، همون کد و همون تست‌بنچ رو یه بار دیگه با عمق ۴ و پهنای ۸ بیت اجرا کردیم --- فقط با سوییچ \lr{-g} موقع لود کردن شبیه‌ساز، بدون یه خط تغییر تو کد. هر هفت مرحله بدون خطا پاس شدن.

\cue{3:55}{سنتز}
تو \lr{Quartus} طرح ۴۰ تا رجیستر و ۱۴ تا پین داد. هر دو عدد رو می‌شه مستقیم از روی کد حساب کرد: ۸ خانه‌ی چهاربیتی حافظه، به‌علاوه‌ی چهار بیت اشاره‌گر، به‌علاوه‌ی چهار بیت \lr{Data\_Out} می‌شه ۴۰؛ و ۸ پین ورودی و ۶ پین خروجی می‌شه ۱۴. اینکه عدد گزارش دقیقاً با محاسبه‌ی دستی می‌خونه، یعنی سنتزکننده همون چیزی رو ساخته که ما در نظر داشتیم.

یه نکته‌ی جالب هم بود: تو \lr{RTL viewer} آرایه‌ی حافظه به‌صورت \lr{SYNC\_RAM} استنتاج شده، ولی گزارش منابع می‌گه صفر بیت حافظه مصرف شده. دلیلش اینه که کلّ حافظه‌ی استک فقط ۳۲ بیته و ابزار ترجیح داده به‌جای اشغال یه بلاک \lr{RAM} اختصاصی، با همون رجیسترهای معمولی بسازتش --- همون ۴۰ رجیستری که تو گزارش می‌بینیم.

آخرش طرح رو روی یه قطعه‌ی مشخّص از \lr{Cyclone II} بردیم و هر ۱۴ پین رو با استاندارد \lr{3.3-V LVTTL} نسبت دادیم. مقایسه‌ی دو قطعه جالب بود: تعداد رجیسترها دقیقاً ثابت موند، چون رجیستر از ساختار طرح میاد نه از فناوری قطعه؛ ولی \lr{Fmax} از ۲۵۲ اومد رو ۲۱۵ مگاهرتز، چون \lr{Cyclone II} فناوری قدیمی‌تریه. تحلیل زمانی هم با قید ۵۰ مگاهرتز حاشیه‌ی مثبت داد.

\cue{4:45}{جمع‌بندی}
درس اصلی این آزمایش فاصله‌ی بین «بیان رفتار» و «بیان ساختار» بود --- کلّ منطق استک تو چند ده خط و بدون هیچ اشاره‌ای به گیت یا فلیپ‌فلاپ نوشته شد و سنتزکننده خودش حافظه و شمارنده و مقایسه‌کننده‌ها رو استنتاج کرد. در کنارش، دقّت رو همون چیزهایی که دستور کار تعریف نکرده بود --- قطبیّت \lr{Empty}، رفتار در برابر سرریز و تهی‌روی، و معنای فرمان هم‌زمان --- تعیین‌کننده‌ی درستی طرح بود. چون آزمایشگاه مجازی بود روی برد فیزیکی پروگرام نکردیم. ممنون، اگه سؤالی هست در خدمتم.

\end{document}
